\chapter{Cơ sở lý thuyết và nghiên cứu liên quan}
\label{chap:background}

Phát hiện vi phạm phân quyền sau khi quyền thay đổi đòi hỏi kết hợp ba nền tảng: kiểm thử REST API dựa trên tìm kiếm, kiểm thử biến hình và kiểm soát truy cập. Chương này trình bày ba nền tảng đó và đối chiếu các công trình liên quan để xác định khoảng trống mà khóa luận nhắm tới.

\section{Kiểm thử REST API dựa trên tìm kiếm}
\label{sec:bg_sbst}

Kiểm thử dựa trên tìm kiếm (search-based software testing) coi việc sinh ca kiểm thử là một bài toán tối ưu hóa. EvoMaster áp dụng cách tiếp cận này cho REST API: công cụ sinh các chuỗi yêu cầu HTTP, rồi dựa vào phản hồi thu được khi thực thi để chọn ca kiểm thử giữ lại \cite{arcuri2019evomaster}. Khóa luận dùng chế độ hộp đen, trong đó EvoMaster chỉ đọc đặc tả OpenAPI và gửi yêu cầu HTTP tới ứng dụng đang chạy \cite{arcuri2024evomastertool}. Thuật toán tìm kiếm mặc định của EvoMaster 6.1.1 tùy theo chế độ chạy. Chế độ hộp trắng dùng Many Independent Objective (MIO), thuật toán chọn và biến đổi ca kiểm thử theo từng mục tiêu \cite{arcuri2017mio}. Chế độ hộp đen dùng SMARTS (Smart Sampling): mỗi vòng lặp, thuật toán lấy mẫu một ca kiểm thử mới và không biến đổi các ca đã có \cite{evomaster611}. Với xác suất 0,95, ca mới được dựng theo chiến lược có cấu trúc, chẳng hạn gọi \texttt{POST} tạo một tài nguyên rồi gọi \texttt{DELETE} lên chính tài nguyên đó; các ca còn lại là chuỗi lời gọi ngẫu nhiên. Các lượt chạy trong khóa luận giữ thuật toán mặc định, nên đều dùng SMARTS. Kho lưu trữ ca kiểm thử giữ những ca đại diện cho các mục tiêu đã đạt; bộ ca kiểm thử cuối và các pha sau đều lấy đầu vào từ kho này.

Pha kiểm thử bảo mật của EvoMaster chạy sau pha tìm kiếm \cite{sahin2024bola}. Pha này dùng các ca kiểm thử và kết quả thực thi trong kho lưu trữ để dựng thêm các ca kiểm thử nhắm vào lỗi bảo mật, không cần thay đổi thuật toán tìm kiếm. Hình~\ref{fig:evomaster_pipeline} biểu diễn chu trình kiểm thử và vị trí pha kiểm thử bảo mật.
\begin{figure}[htbp]
  \centering
  \begin{tikzpicture}[
  >=Stealth,
  font=\small\linespread{1}\selectfont,
  box/.style={rectangle, draw, rounded corners=2pt, fill=gray!10, align=center,
              text width=2.9cm, minimum height=1.2cm},
  ours/.style={box, very thick, fill=gray!30},
  group/.style={draw, dashed, rounded corners=4pt, inner sep=6pt},
  lbl/.style={font=\footnotesize}
]
  \node[box] (spec)  at (0, 0)     {Đặc tả OpenAPI};
  \node[box] (alg)   at (4.1, 0)   {Thuật toán SMARTS};
  \node[box] (exec)  at (8.2, 0)   {Bộ thực thi HTTP};
  \node[box] (sut)   at (12.3, 0)  {Ứng dụng mục tiêu};
  \node[box] (store) at (4.1, -2.0) {Kho lưu trữ ca kiểm thử};

  \node[box]  (gen)   at (4.1, -5.2)  {Bước sinh phản hồi 403 và các oracle có sẵn};
  \node[ours] (build) at (8.2, -5.2)  {Bộ dựng kịch bản\\MR-A1\\{\footnotesize\itshape khóa luận bổ sung}};
  \node[ours] (judge) at (12.3, -5.2) {Bộ phán quyết\\Oracle 970\\{\footnotesize\itshape khóa luận bổ sung}};
  \node[box]  (out)   at (12.3, -7.6) {Bộ ca kiểm thử và danh sách lỗi};

  \begin{scope}[on background layer]
    \node[group, fit=(spec) (alg) (exec) (store)] (g1) {};
    \node[group, fit=(gen) (build) (judge)] (g2) {};
  \end{scope}
  \node[lbl, font=\footnotesize\bfseries, anchor=south west, inner sep=2pt] at (g1.north west) {Pha tìm kiếm};
  \node[lbl, font=\footnotesize\bfseries, anchor=north west, inner sep=3pt] at (g2.south west) {Pha kiểm thử bảo mật};

  \draw[->] (spec) -- (alg);
  \draw[<->] (alg) -- (exec);
  \draw[<->] (exec) -- (sut);
  \draw[->] (alg) -- node[lbl, right] {lưu ca kiểm thử} (store);
  \draw[->] (store) -- (gen);
  \draw[->] (gen) -- (build);
  \draw[->] (build) -- (judge);
  \draw[->, dashed] (build.north) -- ++(0, 0.85) -| node[lbl, pos=0.25, below] {gửi bộ ba HTTP} (sut.south);
  \draw[->] (judge) -- (out);
\end{tikzpicture}

  \caption{Chu trình kiểm thử EvoMaster với pha kiểm thử bảo mật}
  \label{fig:evomaster_pipeline}
\end{figure}

\section{Kiểm thử biến hình}
\label{sec:bg_metamorphic}

Oracle kiểm thử (test oracle) là cơ chế đối chiếu kết quả thực thi với hành vi kỳ vọng hoặc một tiêu chí kiểm tra xác định trước \cite{barr2015oracle}. Khi không thể xây dựng một oracle đầy đủ hoặc chi phí xây dựng quá lớn, quá trình kiểm thử gặp bài toán thiếu oracle (test oracle problem).
Kiểm thử biến hình (metamorphic testing) hỗ trợ kiểm thử khi thiếu oracle bằng cách đối chiếu quan hệ kỳ vọng giữa những lần thực thi có liên quan \cite{chen1998metamorphic}. Một quan hệ biến hình (metamorphic relation) phát biểu mối liên hệ kỳ vọng giữa đầu vào và đầu ra của các lần thực thi đó. Khi các quan sát thỏa điều kiện áp dụng nhưng vi phạm quan hệ, kết quả là bằng chứng về hành vi không phù hợp với thuộc tính được kiểm tra \cite{chen2018metamorphic}.

\section{Lỗ hổng kiểm soát truy cập trong REST API}
\label{sec:bg_access_control}

Kiểm soát truy cập (access control) là cơ chế xác định một chủ thể được phép thực hiện những thao tác nào trên từng tài nguyên của hệ thống. Lỗ hổng phân quyền ở mức đối tượng (Broken Object Level Authorization, BOLA) phát sinh khi API không kiểm tra đủ quyền của chủ thể đối với đối tượng được yêu cầu. Lỗ hổng phân quyền ở mức chức năng (Broken Function Level Authorization, BFLA) phát sinh khi API không kiểm tra đủ quyền gọi một chức năng. Cả hai lớp lỗi nằm trong OWASP API Security Top 10 \cite{owasp2023api}. Lớp lỗi khóa luận xét gần nhất với mục API5 (BFLA) của danh sách này: sau khi bị hạ quyền, tài khoản đã là người dùng thường nhưng vẫn gọi được chức năng của quản trị viên. Điểm khác là lỗi chỉ lộ ra với thông tin xác thực cấp trước lần hạ quyền.

Xác thực (authentication) xác định người gọi là ai, còn phân quyền (authorization) quyết định người gọi được làm gì. Sau khi đăng nhập, máy chủ cấp cho client một thông tin xác thực để gửi kèm các yêu cầu sau. Các ứng dụng trong khóa luận dùng một trong hai cách. Với token JWT, máy chủ ghi tên tài khoản và các thuộc tính như vai trò vào phần claim của token rồi ký, nên ở các yêu cầu sau máy chủ có thể kiểm chữ ký mà không tra cơ sở dữ liệu \cite{rfc7519}. Với phiên phía máy chủ, client giữ một cookie chứa mã phiên, còn dữ liệu của phiên nằm trên máy chủ. Ở cả hai cách, nếu máy chủ quyết định quyền theo vai trò đã chép vào token hay phiên lúc đăng nhập, lần đổi vai trò sau đó chưa có hiệu lực với thông tin xác thực đã cấp cho tới khi nó hết hạn. Có hai hướng khắc phục: tra lại vai trò hiện tại ở mỗi yêu cầu, hoặc thu hồi các thông tin xác thực cấp trước lần đổi vai trò.

Các oracle BOLA và BFLA hiện có đối chiếu phản hồi của các tài khoản có quyền khác nhau. Khác với các oracle này, khóa luận xét trường hợp cùng một chủ thể vẫn được chấp nhận sau khi quyền cần thiết đã bị thu hồi. Hiện tượng này được gọi là lỗ hổng phân quyền có trạng thái (Stateful Broken Access Control). Với hệ thống có chính sách thu hồi quyền có hiệu lực ngay, các yêu cầu tiếp theo phụ thuộc vào quyền vừa bị thu hồi phải bị từ chối. Chuẩn kiểm soát truy cập theo vai trò RBAC mô tả các lệnh thu hồi như hủy gán người dùng khỏi vai trò và thu hồi quyền khỏi vai trò \cite{incits359draft}.

Thay đổi vai trò có áp dụng ngay cho phiên đang hoạt động hay không tùy thuộc vào chính sách và cơ chế thực thi của từng hệ thống. Nếu chính sách yêu cầu chấm dứt hoặc cập nhật phiên mà máy chủ vẫn chấp nhận thông tin xác thực cũ sau khi thu hồi quyền, hành vi này có thể thuộc điểm yếu CWE-1464 (Mục~\ref{sec:intro_problem}).

\section{Các công trình nghiên cứu liên quan}
\label{sec:bg_related_work}

RESTler là một công cụ kiểm thử REST API khác, chủ yếu nhắm tới độ phủ hành vi và lỗi chức năng. Công cụ này sinh các chuỗi lời gọi có trạng thái, kèm một số bộ kiểm tra lỗi \cite{atlidakis2019restler}. Với lỗi phân quyền, Şahin, Zhang và Arcuri bổ sung vào EvoMaster một nhóm oracle an toàn dưới dạng pha hậu xử lý, so sánh nhiều tài khoản tại cùng một thời điểm để phát hiện lỗi BOLA và BFLA \cite{sahin2024bola, arcuri2025issre}. Nhóm oracle này có sẵn trong EvoMaster 6.1.1 và là phương án đối sánh (baseline) của khóa luận.

Ngoài EvoMaster, BACScan phát hiện lỗ hổng kiểm soát truy cập trong ứng dụng web ở chế độ hộp đen; oracle của công cụ dựa vào phản hồi của thao tác để xác định một thao tác đọc hoặc sửa trái phép đã xảy ra hay chưa \cite{liu2025bacscan}. GraphQLer kiểm thử API GraphQL bằng các chuỗi lời gọi nhiều bước: phát lại yêu cầu của một người dùng bằng thông tin xác thực của người dùng khác để tìm lỗi tham chiếu trực tiếp đối tượng không an toàn (IDOR), và chuỗi tạo, xóa, đọc để tìm lỗi dùng tài nguyên sau khi xóa \cite{graphqler2025}. Chuỗi tạo, xóa, đọc có yếu tố thời gian ở phía tài nguyên. Khóa luận xét yếu tố thời gian ở phía quyền: quyền bị thu hồi trong khi thông tin xác thực cấp trước đó vẫn được dùng.

Trong EvoMaster, chưa có oracle nào kết hợp đồng thời ba cơ chế của MR-A1: tự động dựng thao tác thu hồi quyền, dùng lại token cấp trước khi thu hồi, và phán quyết dựa trên quan sát trước và sau khi thay đổi quyền. Nhận định này dựa trên mã nguồn upstream của EvoMaster (bản do nhóm phát triển EvoMaster duy trì) tại commit \texttt{ab983ba338}; Mục~\ref{sec:eval_baseline} kiểm thêm bằng cách chạy bản upstream này trên reports-api. Chương~\ref{chap:method} định nghĩa quan hệ MR-A1 và cách cài đặt quan hệ này trong EvoMaster.
